%---------------------------------------------------------------------
% Quick-reference sheet and the recurring mistakes.
%---------------------------------------------------------------------

\section*{فوری دہرائی}
\markboth{فوری دہرائی}{}

\begin{nukta}[سنین --- تاریخ کے سوال انہی پر جیتے جاتے ہیں]
\begin{center}
\begin{tabular}{@{}p{6.0cm} p{2.6cm} p{4.6cm}@{}}
\toprule
\textbf{واقعہ} & \textbf{ہجری/نبوی} & \textbf{عیسوی} \\
\midrule
پہلی وحی (بعثت)          & \textenglish{--}       & \textenglish{610}ء \\
ہجرتِ حبشہ (پہلی)         & \textenglish{5} نبوی    & \textenglish{615}ء \\
حضرت حمزہؓ و عمرؓ کا اسلام & \textenglish{6} نبوی    & \textenglish{616}ء \\
ہجرتِ مدینہ               & \textenglish{--}       & \textenglish{622}ء \\
غزوہ بدر                  & \textenglish{2}ھ       & \textenglish{624}ء \\
غزوہ اُحد                 & \textenglish{3}ھ       & \textenglish{625}ء \\
غزوہ خندق                 & \textenglish{5}ھ       & \textenglish{627}ء \\
صلح حدیبیہ                & \textenglish{6}ھ       & \textenglish{628}ء \\
غزوہ خیبر                 & \textenglish{7}ھ       & \textenglish{628}ء \\
\textbf{فتح مکہ}          & \textbf{\textenglish{8}ھ} & \textbf{\textenglish{630}ء} \\
غزوہ حنین                 & \textenglish{8}ھ       & \textenglish{630}ء \\
غزوہ تبوک                 & \textenglish{9}ھ       & \textenglish{630}ء \\
\bottomrule
\end{tabular}
\end{center}
\end{nukta}

\medskip
\subsection*{تین اہم فرق}

\begin{center}
\begin{tabular}{@{}p{3.6cm} p{10.8cm}@{}}
\toprule
\textbf{غزوہ / سریہ} & نبی \saw خود شریک ہوں / نبی \saw کسی صحابیؓ کو بھیج
  دیں \\
\textbf{حدیث / سنت} & نبی \saw کا قول / نبی \saw کا مستقل عمل و طرزِ زندگی
  (وسیع تر) \\
\textbf{ایثار / اسراف} & اپنی ضرورت پر دوسرے کو ترجیح / حد سے تجاوز کر کے
  فضول خرچی \\
\bottomrule
\end{tabular}
\end{center}

\medskip
\subsection*{چار عام غلطیاں}

\begin{enumerate}\setlength{\itemsep}{3pt}
  \item \textbf{غزوے میں محض کہانی لکھنا۔} اسباب، واقعات اور نتائج تینوں
        \emph{الگ عنوانوں} سے لکھیے --- محض قصہ سنانا آدھے نمبر دیتا ہے۔
  \item \textbf{ہجری اور عیسوی سن نہ لکھنا۔} تاریخ کے سوال میں سن ہی سب
        سے سستا نمبر ہے۔
  \item \textbf{کامیابی کے نمبر چالیس سمجھنا۔} بہار \textenglish{2022} کے
        پرچے میں یہ \textbf{پچاس} ہیں۔
  \item \textbf{''بحیثیت'' سوالات میں محض تعریف لکھنا۔} ممتحن ٹھوس مثالیں
        اور واقعات مانگتا ہے --- ہر پہلو کے لیے کم از کم دو عملی مثالیں
        لازمی لکھیے۔
\end{enumerate}

\medskip
\noindent
\textbf{\color{bmblue}اگر وقت بہت کم ہو۔}\ \
وہ موضوعات جو \textbf{تین یا چار} پرچوں میں آئے، پہلے تیار کیجیے: سیرت کے
مآخذ (سوال \textenglish{1})، نبی \saw کی کثیر الجہتی شخصیت اور معاشی
تعلیمات (سوال \textenglish{13--16})، اور فتح مکہ (سوال \textenglish{11})۔
ماخذ کی تفصیل کتابچے کے آغاز میں درج ہے۔
